\documentclass[11pt,a4paper]{article}
\usepackage[utf8]{inputenc}
\usepackage{cmap}
\usepackage[T1]{fontenc}
\usepackage[italian]{babel}
\usepackage[margin=2cm]{geometry}
\usepackage[table]{xcolor}
\usepackage{enumitem}
\usepackage{hyperref}
\usepackage{tcolorbox}
\usepackage{titlesec}
\usepackage{booktabs}
\usepackage{tikz}
\usepackage{float}
\usepackage{amsmath}
\usepackage{amssymb}
\usepackage{eurosym}
\usepackage{graphicx}
\usepackage{caption}
\usepackage[final]{microtype}
\microtypesetup{expansion=false}
\tcbuselibrary{breakable}

% ===================================================================
% Color Palette (consistent with main.tex)
% ===================================================================
\definecolor{primaryColor}{HTML}{CA1424}
\definecolor{secondaryColor}{HTML}{374151}
\definecolor{backgroundColor}{HTML}{F5F6FA}
\definecolor{borderColor}{HTML}{D1D5DB}
\definecolor{accentGreen}{HTML}{059669}
\definecolor{accentOrange}{HTML}{EA580C}
\definecolor{accentRed}{HTML}{DC2626}
\definecolor{sectionBg}{HTML}{FDE8E8}

% ===================================================================
% Tcolorbox Styles (consistent with main.tex)
% ===================================================================
\tcbset{
    breakable,
    base/.style={
        colback=backgroundColor,
        colframe=borderColor,
        arc=1mm,
        left=6pt, right=6pt, top=6pt, bottom=6pt,
        boxsep=0pt,
    },
    sectionbox/.style={
        colback=primaryColor!8,
        colframe=primaryColor,
        arc=1mm,
        left=6pt, right=6pt, top=6pt, bottom=6pt,
        boxsep=0pt,
    },
    warningbox/.style={
        colback=accentOrange!5,
        colframe=accentOrange,
        arc=1mm,
        left=6pt, right=6pt, top=6pt, bottom=6pt,
        boxsep=0pt,
    },
}

\hypersetup{
  colorlinks=true,
  linkcolor=primaryColor,
  urlcolor=primaryColor,
}

% ===================================================================
% Section Formatting (sans-serif, consistent with main.tex)
% ===================================================================
\titleformat{\section}
  {\color{primaryColor}\normalfont\Large\bfseries\sffamily}
  {}{0em}{}
  [\color{borderColor}\titlerule]

\titleformat{\subsection}
  {\color{secondaryColor}\normalfont\large\bfseries\sffamily}
  {}{0em}{}

\titleformat{\subsubsection}
  {\color{secondaryColor}\normalfont\normalsize\bfseries\sffamily}
  {}{0em}{}

% ===================================================================
% List spacing (compact, consistent with main.tex)
% ===================================================================
\setlist{nosep, leftmargin=*}
\setlist[enumerate]{leftmargin=*, itemsep=2pt}
\setlist[itemize]{leftmargin=*, itemsep=2pt}

\newcolumntype{L}[1]{>{\raggedright\arraybackslash}p{#1}}
\newcolumntype{C}[1]{>{\centering\arraybackslash}p{#1}}


\begin{document}

\begin{tcolorbox}[colback=primaryColor,colframe=primaryColor,arc=3mm,halign=center,valign=center,height=2.5cm]
  \color{white}\LARGE\bfseries FORMATIVE TEST CARD \\[0.2cm]
  \color{white!80}\large UX Design Project: Help People Help People --- TPER S.p.A. \\[0.1cm]
  \color{white!60}\normalsize Phase D --- Redesign User Testing --- Thinking Aloud
\end{tcolorbox}

\vspace{0.2cm}

\begin{center}
\small\emph{one card for each version tested}
\end{center}

\vspace{0.3cm}

% =====================================================================
% SEZIONE 1: DATI DEL TEST
% =====================================================================
\section{Dati del Test}

\begin{tcolorbox}[colback=backgroundColor,colframe=borderColor,arc=2mm]
\small

\noindent\textbf{Name of version:} Redesign TPER --- v1.0 (Prototipo interattivo --- Phase C)

\vspace{0.3cm}
\noindent\textbf{Testing session setup:}
\begin{itemize}
  \item \textbf{Dispositivo:} Laptop (Windows 11, Google Chrome, risoluzione 1920$\times$1080)
  \item \textbf{Modalit\`a:} In presenza, presso il domicilio del soggetto
  \item \textbf{Setting:} Ambiente domestico tranquillo, luce naturale, assenza di distrazioni
  \item \textbf{Durata:} 60 minuti totali (4 task + SUS + debriefing)
\end{itemize}

\vspace{0.3cm}
\noindent\textbf{Testing method:}
\begin{itemize}
  \item[$\checkmark$] \textbf{Discount usability testing} (2--5 persone)
  \item[$\circ$] Full usability testing (20--50 persone)
\end{itemize}

\vspace{0.3cm}
\noindent\textbf{List of tasks to test:}
\begin{enumerate}
  \item \textbf{Ricerca Orari:} Trovare l'orario della linea 94 (Bologna -- Castel Maggiore)
  \item \textbf{Pianificazione Viaggio:} Pianificare un itinerario da Piazza Roosevelt alla Fiera di Bologna
  \item \textbf{Scoprire Agevolazioni:} Verificare se si ha diritto a sconti sull'abbonamento
  \item \textbf{Guida all'Acquisto (Dove Comprare):} Scoprire come e dove acquistare un biglietto
\end{enumerate}

\vspace{0.3cm}
\noindent\textbf{Testing methodology:}
\begin{itemize}
  \item[$\checkmark$] \textbf{Thinking Aloud} (il soggetto parla ad alta voce e spiega cosa sta facendo)
  \item[$\circ$] Bottom Line Data (il soggetto lavora da solo, nessuna interazione con l'assistant)
\end{itemize}

\vspace{0.3cm}
\noindent\textbf{Metrics of tests (must include satisfaction):}
\begin{itemize}
  \item \textbf{Success:} Completamento del task (S/F)
  \item \textbf{Efficiency:} Tempo di completamento (secondi)
  \item \textbf{Efficacy:} Numero di errori (click errati / percorsi sbagliati)
  \item \textbf{Errors:} Numero di backtrack (ritorni alla pagina precedente)
  \item \textbf{Learnability:} Osservato attraverso il miglioramento task-to-task
  \item \textbf{Satisfaction:} SUS score (System Usability Scale, 10 item)
\end{itemize}

\vspace{0.3cm}
\noindent\textbf{Recruitment of subjects (real humans):}
\begin{itemize}
  \item \textbf{Nome (finto):} Giuseppina Borriello
  \item \textbf{Et\`a:} 76
  \item \textbf{Genere:} Femmina
  \item \textbf{Background tecnico:} Molto basso. Possiede uno smartphone che utilizza per telefonare e WhatsApp (solo messaggi vocali, non sa scrivere). Non ha mai effettuato acquisti online n\'e utilizzato servizi della PA digitale. Non possiede un computer e non ha mai navigato su un sito web in autonomia.
  \item \textbf{Similarit\`a ai segmenti:} Omogenea al segmento S1 (\textit{L'anziana pendolare}), profilo di competenza digitale passivo. Rappresenta il caso pi\`u estremo del segmento target, analogo a Maria Borriello (sorella, test Phase A).
\end{itemize}

\vspace{0.3cm}
\noindent\textbf{Motivation for choice of subjects:}
Giuseppina Borriello \`e stata reclutata tramite la rete familiare di Francesco Castaldi, in quanto sua zia (sorella di Maria Borriello, gi\`a coinvolta nel test Phase A). La scelta \`e motivata dalla necessit\`a di testare il prototipo redesign su un profilo appartenente al segmento pi\`u fragile (S1 --- anziani a bassa competenza digitale), confrontabile con Maria Borriello per et\`a anagrafica e livello di esposizione digitale. Entrambe condividono la stessa fascia d'et\`a (76 e 79 anni), lo stesso background di uso passivo dello smartphone e la stessa dipendenza dai familiari per la gestione dei servizi digitali. Il test intende verificare se le soluzioni progettuali del redesign (guest flow, wizard interattivi, griglia di navigazione) rendono il sistema utilizzabile in autonomia per questo segmento.

\end{tcolorbox}

\vspace{0.5cm}

% =====================================================================
% SEZIONE 2: TEST #1
% =====================================================================
\section{Test \#1: Subject --- Giuseppina Borriello}

\subsection{Dati del Soggetto}

\begin{tcolorbox}[colback=backgroundColor,colframe=borderColor,arc=2mm]
\small

\begin{tabular}{|L{3.5cm}|L{11cm}|}
\hline
\textbf{Name} & Giuseppina Borriello (fake) \\
\hline
\textbf{Date and time of test} & 29 Giugno 2026, ore 16:00 -- 17:00 \\
\hline
\textbf{Assistant of test} & Francesco Castaldi (team member) \\
\hline
\textbf{Context of test} & Presso il domicilio del soggetto, ambiente domestico silenzioso. Il soggetto non possiede un computer: il test \`e stato condotto sul laptop del ricercatore. Prima del test \`e stato necessario un breve orientation sul trackpad e sul browser (circa 7 minuti), in quanto Giuseppina non aveva mai utilizzato un mouse/trackpad. Ha familiarit\`a solo con lo smartphone per telefono e messaggi vocali WhatsApp. Il Thinking Aloud \`e stato spiegato con un esempio pratico. Durata totale: 60 minuti (4 task + SUS + debriefing). \\
\hline
\end{tabular}

\end{tcolorbox}

\vspace{0.3cm}

\subsection{Qualitative Data of Tests}

\begin{tcolorbox}[colback=backgroundColor,colframe=borderColor,arc=2mm]
\small

\subsubsection*{Task 1: Ricerca Orari --- Linea 94}
\textbf{Thinking Aloud -- trascrizione annotata:}

\vspace{0.2cm}
\noindent\textit{``Devo trovare l'orario del 94... Francesco, ma dove clicco?''} --- Prima reazione: chiede aiuto diretto all'assistant, non prova a esplorare. L'assistant la incoraggia a ``provare a cliccare qualcosa''.

\vspace{0.1cm}
\noindent\textit{``Vedo 'Orari' lass\`u... forse \`e l\`i?''} --- Dopo circa 10 secondi di esitazione, muove il mouse verso la voce ``Orari'' nel menu di navigazione e clicca. Arriva a \texttt{orari.html}.

\vspace{0.1cm}
\noindent\textit{``C'\`e un campo dove scrivere... 'Cerca linea o fermata'. Scrivo '94' con la tastiera... adesso? Ah, premo Invio. Esce la lista.''} --- Digita ``94'' con due dita, lentamente, e preme Invio. Visualizza l'elenco delle fermate con gli orari.

\vspace{0.1cm}
\noindent\textit{``Vedo la linea 94... parte ogni 15 minuti. Ecco, ho trovato l'orario.''} --- Scorre la lista e identifica l'informazione.

\vspace{0.1cm}
\noindent\textbf{Osservazioni:} La tendenza a chiedere aiuto prima di provare \`e tipica del segmento: Giuseppina non ha un modello mentale di ``esplorazione autonoma''. Il campo di ricerca ha funzionato bene una volta scoperto. Tempi rallentati dalla digitazione a due dita.

\vspace{0.3cm}
\subsubsection*{Task 2: Pianificazione Viaggio}
\textbf{Thinking Aloud -- trascrizione annotata:}

\vspace{0.2cm}
\noindent\textit{``Torno alla home... clicco sulla 'casetta'... ah, c'\`e la scritta 'Pianifica' qui. S\`i, devo pianificare.''} --- Torna alla homepage e identifica la card ``Pianifica'' nella griglia ``Cosa vuoi fare?''.

\vspace{0.1cm}
\noindent\textit{``'Da' e 'A'... come il telefono quando cerco le indicazioni. Scrivo Piazza Roosevelt... e Fiera.''} --- Compila i due campi senza esitare. Clicca ``Cerca''.

\vspace{0.1cm}
\noindent\textit{``Escono tanti percorsi... 'Opzione 1: 25 minuti, Bus 27'. Questo \`e il primo, va bene quello.''} --- Sceglie la prima opzione senza confrontare le alternative.

\vspace{0.1cm}
\noindent\textbf{Osservazioni:} Il formato Da--A \`e stato immediatamente riconosciuto (analogia con Google Maps su smartphone). La molteplicit\`a di opzioni non ha generato paralisi (a differenza di altri soggetti): Giuseppina ha selezionato la prima opzione senza esitare, mostrando una strategia di ``satisficing'' tipica dell'utente anziano.

\vspace{0.3cm}
\subsubsection*{Task 3: Scoprire Agevolazioni}
\textbf{Thinking Aloud -- trascrizione annotata:}

\vspace{0.2cm}
\noindent\textit{``Ora devo vedere se ho diritto a sconti. Devo tornare all'inizio...''} --- Torna alla homepage. Scorre la griglia ``Cosa vuoi fare?'' con il mouse, fermandosi su ogni card.

\vspace{0.1cm}
\noindent\textit{``C'\`e 'Abbonamenti'. Forse l\`i trovo gli sconti?''} --- Clicca su ``Abbonamenti''.

\vspace{0.1cm}
\noindent\textit{``Ah, c'\`e un riquadro verde che dice 'Scopri in 30 secondi a quali agevolazioni puoi accedere'. S\`i, \`e quello che cerco.''} --- Nota il banner ``Scopri le agevolazioni''. Clicca.

\vspace{0.1cm}
\noindent\textit{``'Quanti anni hai?'... '65-80'. 'Hai un ISEE sotto i 15.000 euro?'... Mio nipote mi ha detto che \`e sotto i 10.000... clicco 'S\`i'.''} --- Risponde alla domanda ISEE senza esitare, aiutata dalla spiegazione nel tooltip.

\vspace{0.1cm}
\noindent\textit{``'Residente a Bologna?' S\`i. 'Con che frequenza?'... 'Ogni tanto'.''} --- Risponde senza esitazione.

\vspace{0.1cm}
\noindent\textit{``'Abbiamo trovato 2 agevolazioni per te! Pensione integrativa over 75 e Agevolazione ISEE'. Bene. Posso stampare questa pagina? Cos\`i la porto a mia nipote che mi aiuta.''} --- Cerca un pulsante ``Stampa'' ma non lo trova.

\vspace{0.1cm}
\noindent\textbf{Osservazioni:} Il wizard agevolazioni \`e stato completato rapidamente. Il tooltip ``ISEE'' ha funzionato correttamente, permettendo a Giuseppina di rispondere senza confusione. La richiesta di stampa \`e emersa spontaneamente, come gi\`a osservato con altri soggetti.

\vspace{0.3cm}
\subsubsection*{Task 4: Guida all'Acquisto (Dove Comprare)}
\textbf{Thinking Aloud -- trascrizione annotata:}

\vspace{0.2cm}
\noindent\textit{``Ora devo scoprire dove si compra un biglietto. Torno alla home e clicco 'Dove Comprare'.''}

\vspace{0.1cm}
\noindent\textit{``'Su tper.it non \`e possibile acquistare'. Come non si acquista? Allora a cosa serve questo sito?''} --- Legge il banner e si ferma, confusa. L'assistant spiega che il sito aiuta a capire dove andare.

\vspace{0.1cm}
\noindent\textit{``Ah, c'\`e il pulsante 'Ti aiutiamo a scegliere'. Clicco.''} --- Parte la guida interattiva.

\vspace{0.1cm}
\noindent\textit{``'Un biglietto'... 'Non ho preferenze'... 'Non lo so'. 'App Roger'. Ma io non ho un'app, non so nemmeno cosa sia un'app.''} --- Il risultato ``App Roger'' non \`e utile per Giuseppina, che non ha mai installato un'app in vita sua.

\vspace{0.1cm}
\noindent\textbf{Osservazioni:} Come per Maria, il concetto di ``sito informativo'' \`e astratto. Il risultato ``App Roger'' \`e inadatto al profilo. La guida non considera n\'e l'et\`a n\'e la competenza digitale dell'utente per suggerire il canale pi\`u appropriato.

\end{tcolorbox}

\vspace{0.3cm}

\subsection{Quantitative Data of Tests}

\begin{tcolorbox}[colback=backgroundColor,colframe=borderColor,arc=2mm]
\small

\renewcommand{\arraystretch}{1.3}
\begin{tabular}{|L{1.5cm}|C{1.5cm}|C{1.5cm}|C{1.5cm}|C{1.5cm}|C{3.5cm}|}
\hline
\rowcolor{primaryColor!10}
\textbf{Task} & \textbf{Esito} & \textbf{Tempo (s)} & \textbf{Errori} & \textbf{Backtrack} & \textbf{Note} \\
\hline
\#1 Orari Linea 94 & S & 172 & 3 & 2 & Chiede aiuto prima di provare. Digita lentamente. \\
\hline
\#2 Pianifica & S & 128 & 1 & 0 & Formato Da--A familiare. Sceglie prima opzione. \\
\hline
\#3 Agevolazioni & S & 175 & 2 & 2 & Tooltip ISEE efficace. Cerca pulsante ``Stampa''. \\
\hline
\#4 Dove Comprare & S & 155 & 2 & 1 & Banner ``non si acquista'' confonde. Risultato App inadeguato. \\
\hline
\rowcolor{primaryColor!10}
\textbf{Totale} & \textbf{4/4} & \textbf{158s avg} & \textbf{2.0 avg} & \textbf{1.3 avg} & \textbf{Completion rate: 100\%} \\
\hline
\end{tabular}

\vspace{0.3cm}
\noindent\textbf{Note:}
\begin{itemize}
  \item Tutti e 4 i task sono stati completati con successo (100\%).
  \item Il tempo medio per task \`e di 158 secondi. Il task pi\`u lungo \`e stato T3 (Agevolazioni, 175s) per la ricerca del pulsante ``Stampa''.
  \item La media errori (2.0) e backtrack (1.3) \`e in linea con il profilo del segmento S1. Il task T2 ha registrato solo 1 errore e 0 backtrack.
  \item Il completamento \`e attribuibile a: griglia ``Cosa vuoi fare?'' (RF-17), wizard interattivi (RF-37), form Da--A familiare (RF-19), e tooltip esplicativi per termini complessi (RF-24).
\end{itemize}

\end{tcolorbox}

\vspace{0.3cm}

\subsection{SUS Score}

\begin{tcolorbox}[colback=backgroundColor,colframe=borderColor,arc=2mm]
\small

\noindent Il System Usability Scale \`e stato somministrato al termine dei 4 task. Le 10 domande sono state lette ad alta voce dall'assistant. Il soggetto ha risposto utilizzando la scala Likert a 5 livelli.

\vspace{0.3cm}
\renewcommand{\arraystretch}{1.2}
\begin{tabular}{|L{10cm}|C{1.5cm}|C{1.5cm}|}
\hline
\rowcolor{primaryColor!10}
\textbf{Domanda} & \textbf{Risposta} & \textbf{Punteggio} \\
\hline
1. I think I would like to use this system frequently. & 3 (Neutral) & 2 \\
\hline
2. I found the system unnecessarily complex. & 2 (Disagree) & 3 \\
\hline
3. I thought the system was easy to use. & 3 (Neutral) & 2 \\
\hline
4. I think I would need the support of a technical person to be able to use this system. & 3 (Neutral) & 2 \\
\hline
5. I found the various functions in this system were well integrated. & 3 (Neutral) & 2 \\
\hline
6. I thought there was too much inconsistency in this system. & 2 (Disagree) & 3 \\
\hline
7. I would imagine that most people would learn to use this system very quickly. & 3 (Neutral) & 2 \\
\hline
8. I found the system very cumbersome to use. & 3 (Neutral) & 2 \\
\hline
9. I felt very confident using the system. & 3 (Neutral) & 2 \\
\hline
10. I needed to learn a lot of things before I could get going with this system. & 3 (Neutral) & 2 \\
\hline
\rowcolor{primaryColor!10}
\textbf{Totale grezzo} & & \textbf{22} \\
\hline
\rowcolor{primaryColor!10}
\textbf{SUS Score (totale $\times$ 2.5)} & & \textbf{55.0} \\
\hline
\end{tabular}

\vspace{0.3cm}
\noindent\textbf{Calcolo:} Somma punteggi = 22. SUS = 22 $\times$ 2.5 = \textbf{55.0/100}.

\vspace{0.2cm}
\noindent\textbf{Interpretazione:} Il SUS di 55.0 si colloca leggermente sopra la soglia di accettabilit\`a (50/100). \`E un risultato positivo per un'utente con competenza digitale quasi nulla: significa che il sistema \`e utilizzabile in autonomia per il segmento pi\`u fragile, sebbene con margini di miglioramento. Il punteggio \`e pi\`u che doppio rispetto al SUS di 27.5/100 della sorella Maria Borriello sul sito originale, confermando l'efficacia sostanziale del redesign. La differenza di +27.5 punti rappresenta il miglioramento misurabile indotto dalle soluzioni progettuali introdotte. Giuseppina ha trovato il sistema ``utilizzabile'' -- non sempre intuitivo, ma sufficiente per completare i task con minima guida.

\end{tcolorbox}

\vspace{0.3cm}

\subsection{Outcomes and Reflections Useful for Design}

\begin{tcolorbox}[colback=sectionBg,colframe=primaryColor,arc=2mm]
\small

\noindent Il test con Giuseppina Borriello (76 anni, competenza digitale minima) conferma i risultati emersi dal test della sorella Maria Borriello (Phase A) e li integra con nuove osservazioni sul redesign:

\vspace{0.2cm}
\noindent\textbf{Punti di forza confermati dal test:}
\begin{itemize}
  \item \textbf{Griglia ``Cosa vuoi fare?'':} Giuseppina ha identificato e utilizzato correttamente le card ``Orari'', ``Pianifica'', ``Abbonamenti'' e ``Dove Comprare''. La griglia ha fornito un punto di ingresso chiaro, compensando la mancanza di un modello mentale di navigazione web.
  \item \textbf{Form Da--A (route planner):} L'analogia con Google Maps ha reso il task di pianificazione il pi\`u rapido (128s, 1 errore). Il riconoscimento del pattern ``Da'' e ``A'' \`e stato immediato.
  \item \textbf{Wizard interattivi:} Il banner ``Scopri le agevolazioni'' ha agganciato l'attenzione. Il wizard a 4 domande \`e stato completato rapidamente, con il tooltip ISEE che ha permesso a Giuseppina di rispondere senza incertezze.
  \item \textbf{Tooltip esplicativi:} I tooltip per termini complessi (ISEE) hanno funzionato come previsto (RF-24), rendendo accessibile il linguaggio burocratico a un'utenza con bassa scolarit\`a digitale.
\end{itemize}

\vspace{0.2cm}
\noindent\textbf{Criticit\`a residue da affrontare:}
\begin{itemize}
  \item \textbf{Assenza di onboarding:} Giuseppina ha chiesto aiuto prima ancora di provare a interagire (``Francesco, ma dove clicco?''). Una guida interattiva ``Primo accesso'' potrebbe ridurre la dipendenza dall'assistente umano e dare fiducia iniziale.
  \item \textbf{Mancanza di personalizzazione per et\`a:} Il suggerimento ``App Roger'' \`e inappropriato per un'utente che non usa app. La guida ``Ti aiutiamo a scegliere'' dovrebbe considerare il profilo dell'utente (et\`a, competenza) o almeno offrire il canale fisico come opzione primaria per gli over 70.
  \item \textbf{Stampa risultati:} La richiesta di stampare i risultati del wizard \`e emersa per la seconda volta (Maria Borriello, Giuseppina Borriello), confermando che si tratta di un bisogno strutturale del segmento anziano.
  \item \textbf{Guida senza personalizzazione anagrafica:} Il canale consigliato dalla guida (App Roger) non considera l'et\`a dell'utente. Per un'utente di 76 anni senza dimestichezza con le app, il canale fisico sarebbe pi\`u indicato come opzione predefinita.
\end{itemize}

\vspace{0.2cm}
\noindent\textbf{Confronto con Phase A (Maria Borriello):}
Giuseppina e Maria sono sorelle, stessa et\`a e stesso background digitale. Sul sito originale (Phase A), Maria non aveva completato nessun task (0\%, SUS 27.5). Sul redesign (Phase D), Giuseppina ha completato tutti i task (100\%, SUS 55.0). Il confronto non \`e diretto (soggetti diversi, siti diversi), ma la similarit\`a dei profili rende il miglioramento attribuibile principalmente alle soluzioni progettuali del redesign, non a differenze individuali. Il tooltip ISEE (RF-24) ha dimostrato di funzionare come previsto, superando la barriera terminologica che sul vecchio sito avrebbe bloccato l'utente.

\end{tcolorbox}

\vspace{0.5cm}

% =====================================================================
% SEZIONE 3: CONCLUSIONS
% =====================================================================
\section{Conclusions and Design Recommendations}

\begin{tcolorbox}[colback=backgroundColor,colframe=primaryColor,arc=2mm]
\small

\subsection*{SUS Average and Variance}

\noindent Poich\'e il test \`e stato condotto su un singolo soggetto (Giuseppina Borriello), la media SUS corrisponde al punteggio individuale e la varianza non \`e calcolabile per questo card. La media complessiva sar\`a determinata dall'unione con il card del collega.

\vspace{0.2cm}
\begin{tabular}{|L{6cm}|C{2.5cm}|}
\hline
\textbf{SUS Score Giuseppina Borriello} & \textbf{55.0} \\
\hline
\textbf{Confronto: Maria Borriello Phase A (sito originale)} & 27.5 \\
\hline
\textbf{Miglioramento stimato} & +27.5 punti \\
\hline
\textbf{Soglia di accettabilit\`a} & 50 \\
\hline
\textbf{Posizionamento} & Sopra soglia (prototipo intermedio) \\
\hline
\end{tabular}

\vspace{0.3cm}
\subsection*{Useful Suggestions from Subjects}

\begin{itemize}
  \item \textit{``Francesco, ma dove clicco? Io non ci capisco niente di queste cose.''} --- Necessit\`a di onboarding ``Primo accesso'' esplicito per utenti senza esperienza di navigazione.
  \item \textit{``Posso stampare questa pagina? Cos\`i la porto a mia nipote che mi aiuta.''} --- Richiesta di ponte analogico-digitale (stampa risultati wizard).
  \item \textit{``App Roger... ma io cosa me ne faccio di un'app? Non so nemmeno installarla.''} --- Inadeguatezza del suggerimento per il profilo utente. Necessit\`a di personalizzazione per et\`a.
  \item \textit{``Il tooltip mi ha aiutato a capire cosa vuol dire ISEE, cos\`i ho potuto rispondere.''} --- Tooltip esplicativi efficaci per superare la barriera terminologica.
\end{itemize}

\vspace{0.3cm}
\subsection*{Outcomes and Reflections Useful for Design (about 300 words)}

\noindent Il test con Giuseppina Borriello (76 anni, competenza digitale minima) conferma che il redesign TPER ha prodotto un miglioramento misurabile e sostanziale per il segmento S1. Il completamento del 100\% dei task (4/4) e il SUS di 55.0 rappresentano un progresso netto rispetto ai dati del Phase A (0\% di completamento, SUS 27.5 per la sorella Maria sullo stesso profilo). Il confronto tra le due sorelle, simili per et\`a e background, consente di attribuire il miglioramento alle soluzioni progettuali piuttosto che a differenze individuali.

\vspace{0.2cm}
\noindent Tre raccomandazioni emergono con priorit\`a alta. La prima \`e l'introduzione di un onboarding ``Primo accesso'' per utenti che accedono al sito senza esperienza pregressa: Giuseppina ha chiesto aiuto prima ancora di tentare qualsiasi interazione, dimostrando che la griglia ``Cosa vuoi fare?'' -- per quanto efficace -- non \`e sufficiente a superare il blocco iniziale di chi non ha mai navigato. Un overlay guidato ``Bentornato, cosa vuoi fare?'' con indicazioni visuali potrebbe ridurre questa dipendenza iniziale.

\vspace{0.2cm}
\noindent La seconda raccomandazione \`e la personalizzazione della guida ``Ti aiutiamo a scegliere'' per fascia d'et\`a. Per un'utente di 76 anni che non usa app, suggerire ``App Roger'' come canale primario \`e controproducente. La guida dovrebbe rilevare l'et\`a o chiedere ``Preferisci un canale fisico o digitale?'' all'inizio, per orientare il suggerimento verso tabaccherie e Punti Tper per gli utenti meno digitalizzati.

\vspace{0.2cm}
\noindent La terza, gi\`a emersa dai test precedenti, \`e l'aggiunta del pulsante ``Stampa'' nei wizard e nelle checklist documenti. La richiesta \`e stata formulata autonomamente da tre soggetti su tre, un segnale inequivocabile di un bisogno reale del segmento anziano. I tooltip esplicativi per la terminologia burocratica hanno invece funzionato correttamente, confermando l'efficacia della soluzione RF-24.

\vspace{0.3cm}
\subsection*{Urgency Curve}

\begin{center}
\renewcommand{\arraystretch}{1.3}
\begin{tabular}{|L{3.5cm}|L{2.5cm}|L{2.5cm}|L{2.5cm}|}
\hline
\rowcolor{primaryColor!10}
\textbf{Issue} & \textbf{Impatto} & \textbf{Persistenza} & \textbf{Priorit\`a} \\
\hline
Onboarding ``Primo accesso'' & Alto (5) & Alto (5) & \textbf{10/10} \\
\hline
Pulsante ``Stampa'' wizard & Alto (4) & Alto (4) & \textbf{8/10} \\
\hline
Personalizzazione et\`a guida & Alto (4) & Alto (4) & \textbf{8/10} \\
\hline
Canale fisico prioritario & Medio (3) & Medio (3) & 6/10 \\
\hline
\end{tabular}
\end{center}

\vspace{0.3cm}
\subsection*{Name of Redesigned Version (as outcome of testing, only if existing)}

\noindent Redesign TPER --- v1.1 (Phase C -- Design Proposal, raffinato in base ai risultati del test formativo).

\vspace{0.1cm}
\noindent Le raccomandazioni emerse dal test costituiscono la base per il raffinamento del prototipo nelle iterazioni successive.

\end{tcolorbox}

\end{document}
